\documentclass[11pt]{article}
\usepackage[utf8]{inputenc}
\usepackage[T1]{fontenc}
\usepackage{amsmath}
\usepackage{amssymb}
\usepackage{multicol}
\usepackage{geometry}
\usepackage{enumitem}
\usepackage{xcolor}
\usepackage{titlesec}

% Layout
\geometry{a4paper, left=1cm, right=1cm, top=0.5cm, bottom=1.2cm}
\setlength{\columnseprule}{0.4pt}

% Títulos
\titleformat{\section}{\normalfont\Large\bfseries\color{blue}}{}{0em}{}
\titleformat{\subsection}{\normalfont\large\bfseries\color{red}}{}{0em}{}

\title{\textcolor{red}{Segundo Trimestre: Matemática\\Razão e Proporção}}
\author{Professor: Jefferson}
\date{}

\begin{document}

\maketitle
\vspace{-1cm}

\begin{multicols}{2}

\section*{Questões Comentadas}

\subsection*{Questão 1}
A razão entre dois números é $3:5$ e a soma desses números é 64. Determine os números.

\textbf{Comentário:}
\[
\frac{3}{5} = \frac{x}{64-x}
\Rightarrow 5x = 192 - 3x
\Rightarrow x = 24
\]
Logo, os números são 24 e 40.

\textbf{Resposta:} 24 e 40

\subsection*{Questão 2}
A razão entre as idades de Ana e Beatriz é $2:3$. Daqui a 4 anos, a razão será $3:4$. Qual é a idade atual de Ana?

\textbf{Comentário:}
\[
\frac{2x}{3x} \Rightarrow
\frac{2x + 4}{3x + 4} = \frac{3}{4}
\Rightarrow 8x + 16 = 9x + 12
\Rightarrow x = 4
\]
Idade de Ana: $2 \cdot 4 = 8$.

\textbf{Resposta:} 8 anos

\subsection*{Questão 3}
Três grandezas $A$, $B$ e $C$ são diretamente proporcionais. Sabendo que
\[
A = 4,\ B = 6,\ C = 10
\]
e que $A = 10$, determine o novo valor de $C$.

\textbf{Comentário:}
\[
\frac{4}{10} = \frac{10}{x}
\Rightarrow x = 25
\]

\textbf{Resposta:} 25

\subsection*{Questão 4}
Uma máquina A produz 120 peças em 6 horas. Uma máquina B, mais eficiente, produz a mesma quantidade em 4 horas. Qual é a razão entre as eficiências das máquinas?

\textbf{Comentário:}
\[
\text{A}: \frac{120}{6} = 20
\quad
\text{B}: \frac{120}{4} = 30
\]
\[
\text{Razão} = \frac{20}{30} = \frac{2}{3}
\]

\textbf{Resposta:} $2:3$

\subsection*{Questão 5}
Em uma mistura, a razão entre álcool e água é $3:5$. Se a mistura possui 32 litros, quantos litros de álcool há?

\textbf{Comentário:}
\[
\frac{3}{8} = \frac{x}{32}
\Rightarrow x = 12
\]

\textbf{Resposta:} 12 litros

\subsection*{Questão 6}
4 operários constroem um muro em 15 dias trabalhando 6 horas por dia. Quantos dias levarão 6 operários trabalhando 5 horas por dia?

\textbf{Comentário:}
\[
4 \cdot 15 \cdot 6 = 6 \cdot x \cdot 5
\Rightarrow x = 12
\]

\textbf{Resposta:} 12 dias

\subsection*{Questão 7}
Um mapa utiliza a escala 1:250000. Qual a distância real correspondente a 3,2 cm no mapa?

\textbf{Comentário:}
\[
3{,}2 \cdot 250000 = 800000 \text{ cm} = 8 \text{ km}
\]

\textbf{Resposta:} 8 km

\subsection*{Questão 8}
A razão entre o número de alunos de duas turmas é $7:9$. Sabendo que a diferença entre elas é 12 alunos, determine o total de alunos.

\textbf{Comentário:}
\[
9x - 7x = 12 \Rightarrow x = 6
\]
Total:
\[
7x + 9x = 16x = 96
\]

\textbf{Resposta:} 96 alunos

\subsection*{Questão 9}
Um carro consome combustível de forma proporcional à distância percorrida. Se percorre 180 km com 12 litros, quantos litros gastará em 75 km?

\textbf{Comentário:}
\[
\frac{180}{75} = \frac{12}{x}
\Rightarrow x = 5
\]

\textbf{Resposta:} 5 litros

\subsection*{Questão 10}
Dois números estão na razão $5:8$. Se o maior excede o menor em 27 unidades, determine o maior número.

\textbf{Comentário:}
\[
8x - 5x = 27 \Rightarrow x = 9
\Rightarrow 8x = 72
\]

\textbf{Resposta:} 72

\subsection*{Questão 11}
6 máquinas produzem 900 peças em 15 dias. Quantas máquinas são necessárias para produzir 1200 peças em 10 dias?

\textbf{Comentário:}
\[
6 \cdot 15 \cdot 900 = x \cdot 10 \cdot 1200
\Rightarrow x = 6{,}75
\]

\textbf{Resposta:} 7 máquinas (aproximadamente)

\subsection*{Questão 12}
A razão entre dois números é $\frac{3}{4}$. Se o menor é 45, determine o maior.

\textbf{Comentário:}
\[
\frac{3}{4} = \frac{45}{x}
\Rightarrow x = 60
\]

\textbf{Resposta:} 60

\subsection*{Questão 13}
Se 8 torneiras enchem um tanque em 6 horas, em quanto tempo 12 torneiras iguais encherão o mesmo tanque?

\textbf{Comentário:}
\[
8 \cdot 6 = 12 \cdot x
\Rightarrow x = 4
\]

\textbf{Resposta:} 4 horas

\subsection*{Questão 14}
Uma receita usa farinha, açúcar e manteiga na razão $5:3:2$. Se foram usados 600 g de farinha, quantos gramas de açúcar foram usados?

\textbf{Comentário:}
\[
\frac{3}{5} = \frac{x}{600}
\Rightarrow x = 360
\]

\textbf{Resposta:} 360 g

\subsection*{Questão 15}
A razão entre a velocidade de dois atletas é $4:5$. Se o mais rápido percorre 100 m em 20 s, quanto tempo o outro leva?

\textbf{Comentário:}
\[
\frac{4}{5} = \frac{20}{x}
\Rightarrow x = 25
\]

\textbf{Resposta:} 25 s

\subsection*{Questão 16}
Se 3 kg de tinta pintam 45 m², quantos metros quadrados podem ser pintados com 8 kg?

\textbf{Comentário:}
\[
\frac{3}{8} = \frac{45}{x}
\Rightarrow x = 120
\]

\textbf{Resposta:} 120 m²

\subsection*{Questão 17}
Dois números estão na razão $2:7$ e sua soma é 81. Determine o maior número.

\textbf{Comentário:}
\[
2x + 7x = 81
\Rightarrow x = 9
\Rightarrow 7x = 63
\]

\textbf{Resposta:} 63

\subsection*{Questão 18}
Um trabalhador recebe salário proporcional aos dias trabalhados. Se ganhou R\$ 1.260 por 18 dias, quanto ganhará por 25 dias?

\textbf{Comentário:}
\[
\frac{18}{25} = \frac{1260}{x}
\Rightarrow x = 1750
\]

\textbf{Resposta:} R\$ 1.750

\subsection*{Questão 19}
A razão entre as áreas de dois quadrados é $9:16$. Qual é a razão entre seus lados?

\textbf{Comentário:}
\[
\sqrt{\frac{9}{16}} = \frac{3}{4}
\]

\textbf{Resposta:} $3:4$

\subsection*{Questão 20}
Explique por que a regra de três é uma aplicação direta do conceito de proporção.

\textbf{Comentário:}
Ela utiliza a igualdade entre duas razões para encontrar um valor desconhecido.

\textbf{Resposta:} Porque resolve proporções por igualdade de razões.

\end{multicols}

\end{document}

